\subsection{Implications for reported results}
The seven corpora fall into three groups. PhiUSIIL and the PT-Alexa construction are separable by recording conventions, and the contrast persists among homepages alone (Table~\ref{tab:controls}): \texttt{https://www.} in front of every legitimate PhiUSIIL string, a missing scheme in front of every legitimate PT-Alexa string, and no path in any legitimate URL of either. On these corpora raw-string scores say little about phishing detection; canonical results are the more informative numbers, and they are lower by \DRawVsCanonPhiUSIIL{} and \DRawVsCanonPTAlexa{} MCC. Because their legitimate classes contain only homepages, even canonical results there partly measure the homepage contrast, which no transformation of the strings can remove.

Kaggle-MU is a different case. Its character model depends on formatting (a drop of \DRawVsCanonKaggleMU, \KCleanDrop{} on the cleaned subset and \KCleanTwoDrop{} without the bare-\texttt{www.} phishing rows), but the larger problem is its labels: thousands of canonical strings and registered domains with both labels, a standards body as its largest phishing domain, thousands of URLs that PhishStorm labels the other way, and thousands of benign-looking pages labelled phishing. Formatting and source mixing cannot be separated in this collection. It also contains \ISCXInKagglePct\% of ISCX-URL2016. It should not be used as an independent benchmark next to ISCX-URL2016 or PhishStorm. Its low transfer scores are consistent with these label inconsistencies, although we have not tested that they cause them.

Ebbu2017, PhishStorm, ISCX-URL2016 and Hannousse2021 have moderate presentation probes, but a character n-gram model loses less than 0.01 MCC on average when the formatting is removed (a CNN up to \CNNDropHannousseRandom{} under the random split), and grouping by registered domain costs the canonical n-gram model at most \DCanonRandomVsDomainPhishStorm{} MCC there, the lexical model up to \DLexicalRandomVsDomainISCXURL{} and the CNN up to \CNNGroupCostISCXURL{} on ISCX-URL2016. They are the most usable of the seven for URL-only detection, with one caveat: ISCX-URL2016 draws its legitimate class from \RegLegitISCXURL{} registered domains, so any detector that weights the host heavily needs a grouped split there.

For a security operations team the operating point matters more than the balanced score. Within a corpus, a threshold set on validation data for a false-positive rate of 0.1\% holds that rate on the test data and gives the canonical model a precision between \VPrecOneKaggleMUCharCanon{} and \VPrecOneISCXURLCharCanon{} at 1\% prevalence (Table~\ref{tab:precision}), at true-positive rates between \VTprKaggleMUCharCanon{} and \VTprISCXURLCharCanon. Across corpora the same procedure fails: the realised false-positive rate has a median of \XVFprCharCanon\% and exceeds ten times the target in \XVOverTenCharCanon{} of \XVNPairsCharCanon{} pairs, and the median precision at 1\% prevalence would be \XVPrecOneCharCanon, close to the base rate itself. A threshold therefore has to be set on traffic from the deployment, and a benchmark score does not tell an operator where to set it. A filter tuned on PhiUSIIL raw strings would also have learned that \texttt{https://www.} homepages are legitimate, which is the opposite of what a defender wants once phishing sites use HTTPS. These findings do not show that published detectors are weak. They show that reported numbers mix detection skill with corpus-specific conventions in unknown proportions, and that the proportions differ sharply between corpora.

What a deployed detector sees also differs from what these corpora store. A mail gateway or a browser extension sees full URLs with a scheme, often wrapped by a rewriting service or a click tracker that has to be unwrapped first; a DNS resolver, or a proxy without TLS interception, sees only the host name. For the second group the host-view results are the relevant ones, and they are the results most affected by domain leakage. Phishing hosted on legitimate platforms is a further limit of URL-only detection: \texttt{\LargestPhishPhiUSIIL}, \texttt{\LargestPhishPTAlexa} and \texttt{\LargestPhishISCXURL} are the registered domains with the most phishing URLs in PhiUSIIL, PT-Alexa and ISCX-URL2016, so a host-only model cannot separate those URLs from benign pages on the same platform.
